\documentclass[12pt,a4paper]{article}
\usepackage[utf8]{inputenc}
\usepackage[T1]{fontenc}
\usepackage[indonesian]{babel}
\usepackage{geometry}
\usepackage{longtable}
\usepackage{array}
\usepackage{booktabs}
\usepackage{enumitem}
\usepackage[hidelinks,breaklinks]{hyperref}
\geometry{margin=2.5cm}

\begin{document}

\begin{center}
\textbf{\Large RENCANA PEMBELAJARAN SEMESTER (RPS)}\\
\textbf{Berbasis Outcome-Based Education (OBE)}\\[0.5cm]
\textbf{PROGRAM STUDI TEKNIK INFORMATIKA}\\
\textbf{FAKULTAS TEKNOLOGI INFORMASI}\\
\textbf{UNIVERSITAS BALE BANDUNG}\\
\end{center}

\vspace{1cm}

% ============================================================
% 1. IDENTITAS MATA KULIAH
% ============================================================
\section*{1. Identitas Mata Kuliah}
\begin{tabular}{ll}
Nama Program Studi & : Teknik Informatika \\
Nama Mata Kuliah & : Data Science \\
Kode Mata Kuliah & : TIF-xxx \\
Semester & : 5 (Lima) \\
SKS / Bobot Kredit & : 3 SKS (komposisi 60\% teori, 40\% praktek) \\
Jumlah Pertemuan & : 16 (enam belas) pertemuan \\
Dosen Pengampu & : [Nama Dosen, M.Kom.] \\
Tanggal Penyusunan & : [Tanggal] \\
\end{tabular}

\vspace{0.5cm}

% ============================================================
% 2. CPL PRODI YANG DIBEBANKAN PADA MATA KULIAH
% ============================================================
\section*{2. CPL Prodi yang Dibebankan pada Mata Kuliah}

CPL Prodi berikut dipetakan pada capaian mata kuliah Data Science (selaras dengan dokumen RPS):

\begin{itemize}[leftmargin=*]
  \item \textbf{CPL01} --- Kemampuan menganalisis persoalan \textit{computing} yang kompleks serta menerapkan prinsip-prinsip \textit{computing} dan disiplin ilmu relevan lainnya untuk mengidentifikasi solusi, dengan mempertimbangkan wawasan perkembangan ilmu transdisiplin.

  \item \textbf{CPL03} --- Kemampuan merancang dan menganalisis algoritma untuk menyelesaikan permasalahan organisasi secara optimal, serta memilih dan menerapkannya pada bahasa pemrograman tertentu.
\end{itemize}

\vspace{0.5cm}

% ============================================================
% 3. CAPAIAN PEMBELAJARAN MATA KULIAH (CPMK)
% ============================================================
\section*{3. Capaian Pembelajaran Mata Kuliah (CPMK)}

Kemampuan atau kompetensi spesifik yang diharapkan mahasiswa kuasai setelah menyelesaikan mata kuliah (selaras dengan CPL Prodi yang dibebankan):

\begin{itemize}[leftmargin=*]
  \item \textbf{CPMK1} --- (Fokus pemahaman masalah dan wawasan transdisiplin; terkait CPL01) Mahasiswa mampu menganalisis persoalan komputasi dan bisnis yang kompleks dari berbagai domain (transdisiplin), serta merumuskan pendekatan berbasis data (\textit{data-driven}) sebagai solusi strategis untuk organisasi.

  \item \textbf{CPMK2} --- (Fokus eksplorasi data dan pemrograman; terkait CPL01 dan CPL03) Mahasiswa mampu menerapkan prinsip-prinsip komputasi untuk melakukan pengumpulan, prapemrosesan (\textit{data preprocessing}), dan Analisis Data Eksploratif (EDA) menggunakan bahasa pemrograman tertentu (seperti Python atau R) guna mengidentifikasi pola tersembunyi pada data.

  \item \textbf{CPMK3} --- (Fokus perancangan dan pemilihan algoritma; terkait CPL03) Mahasiswa mampu merancang, memilih, dan menerapkan algoritma \textit{Machine Learning} atau analisis statistik yang paling optimal (seperti regresi, klasifikasi, atau klastering) dalam bahasa pemrograman untuk menyelesaikan permasalahan organisasi secara efektif.

  \item \textbf{CPMK4} --- (Fokus evaluasi dan penyampaian solusi; terkait CPL01 dan CPL03) Mahasiswa mampu menganalisis dan mengevaluasi kinerja model/algoritma yang telah dibangun, serta mengkomunikasikan hasil analisis teknis menjadi wawasan (\textit{insight}) bisnis yang mudah dipahami guna mendukung pengambilan keputusan.
\end{itemize}

\vspace{0.5cm}

% ============================================================
% 4. SUB-CPMK / INDIKATOR PENCAPAIAN
% ============================================================
\section*{4. Sub-CPMK / Indikator Pencapaian}

Penjabaran CPMK menjadi indikator yang lebih terukur dan dapat diuji (selaras dengan RPS):

\begin{itemize}[leftmargin=*]
  \item \textbf{Sub-CPMK 1:} Mahasiswa mampu mengidentifikasi permasalahan bisnis dari berbagai domain (transdisiplin) dan menjelaskan siklus hidup data science sebagai solusi \textit{data-driven}.

  \item \textbf{Sub-CPMK 2:} Mahasiswa mampu melakukan pengumpulan data dari berbagai sumber serta memahami struktur dan karakteristik awal data.

  \item \textbf{Sub-CPMK 3:} Mahasiswa mampu menerapkan teknik prapemrosesan data (\textit{data cleaning}) untuk menangani \textit{missing values}, anomali, dan duplikasi menggunakan Python/R.

  \item \textbf{Sub-CPMK 4:} Mahasiswa mampu melakukan transformasi data dan rekayasa fitur (\textit{feature engineering}) agar data optimal untuk proses pemodelan.

  \item \textbf{Sub-CPMK 5:} Mahasiswa mampu merancang dan mengeksekusi Analisis Data Eksploratif (EDA) serta memvisualisasikannya untuk menemukan pola tersembunyi.

  \item \textbf{Sub-CPMK 6:} Mahasiswa mampu merumuskan hipotesis dan memilih algoritma yang relevan berdasarkan jenis masalah (\textit{Supervised} vs \textit{Unsupervised Learning}).

  \item \textbf{Sub-CPMK 7:} Mahasiswa mampu merancang dan mengimplementasikan algoritma regresi untuk menyelesaikan permasalahan prediksi kuantitatif pada organisasi.

  \item \textbf{Sub-CPMK 8:} Mahasiswa mampu mengimplementasikan algoritma klasifikasi dasar untuk memprediksi kelas atau kategori dari suatu permasalahan bisnis (misal: \textit{churn prediction}).

  \item \textbf{Sub-CPMK 9:} Mahasiswa mampu menerapkan algoritma klasifikasi tingkat lanjut (\textit{ensemble}) untuk meningkatkan keandalan dan akurasi prediksi model.

  \item \textbf{Sub-CPMK 10:} Mahasiswa mampu merancang algoritma klastering untuk mengidentifikasi segmentasi data yang tidak berlabel (misal: segmentasi pelanggan).

  \item \textbf{Sub-CPMK 11:} Mahasiswa mampu menganalisis dan menghitung metrik evaluasi kinerja model/algoritma yang telah dibangun secara presisi.

  \item \textbf{Sub-CPMK 12:} Mahasiswa mampu mengoptimalkan kinerja algoritma melalui teknik validasi dan penyesuaian parameter (\textit{hyperparameter tuning}).

  \item \textbf{Sub-CPMK 13:} Mahasiswa mampu mengkomunikasikan hasil evaluasi teknis algoritma menjadi narasi dan visualisasi (\textit{insight}) yang relevan untuk pengambilan keputusan bisnis.

  \item \textbf{Sub-CPMK 14:} Mahasiswa mampu merancang dan mengeksekusi proyek data science secara utuh (\textit{end-to-end pipeline}) sebagai solusi strategis untuk studi kasus nyata.
\end{itemize}

\vspace{0.5cm}

% ============================================================
% 5. MATERI PEMBELAJARAN (BAHAN KAJIAN)
% ============================================================
\section*{5. Materi Pembelajaran (Bahan Kajian)}

Rincian materi per pertemuan beserta kaitannya dengan Sub-CPMK dan CPMK (selaras dengan RPS dan 16 pertemuan):

\vspace{0.35cm}
{\small\setlength{\tabcolsep}{3.5pt}%
\begin{longtable}{|>{\raggedright\arraybackslash}p{2.05cm}|>{\raggedright\arraybackslash}p{4.55cm}|>{\raggedright\arraybackslash}p{4.55cm}|>{\raggedright\arraybackslash}p{2.95cm}|}
\hline
\textbf{Pertemuan} & \textbf{Materi pembelajaran} & \textbf{Sub-CPMK} & \textbf{Kaitan CPMK} \\
\hline
\endfirsthead
\hline
\textbf{Pertemuan} & \textbf{Materi pembelajaran} & \textbf{Sub-CPMK} & \textbf{Kaitan CPMK} \\
\hline
\endhead
\hline
\endfoot

\textbf{Pertemuan 1} &
Pengantar Data Science, Siklus Hidup Data (CRISP-DM), Pemahaman Bisnis (\textit{Business Understanding}). &
\textbf{Sub-CPMK 1:} Mahasiswa mampu mengidentifikasi permasalahan bisnis dari berbagai domain dan menjelaskan siklus hidup data science. &
\textbf{CPMK1} (fokus pemahaman masalah dan wawasan transdisiplin). \\
\hline

\textbf{Pertemuan 2} &
Teknik Pengumpulan Data (\textit{Web Scraping}, API, Database), Struktur Data, \textit{Data Understanding}. &
\textbf{Sub-CPMK 2:} Mahasiswa mampu melakukan pengumpulan data dari berbagai sumber serta memahami struktur dan karakteristik awal data. &
\textbf{CPMK2} (fokus eksplorasi data dan pemrograman). \\
\hline

\textbf{Pertemuan 3} &
\textit{Data Preprocessing} 1: \textit{Missing values imputation}, \textit{Outlier handling}, \textit{Data deduplication}. &
\textbf{Sub-CPMK 3:} Mahasiswa mampu menerapkan teknik prapemrosesan data (\textit{data cleaning}) untuk menangani \textit{missing values}, anomali, dan duplikasi menggunakan Python/R. &
\textbf{CPMK2} (fokus eksplorasi data dan pemrograman). \\
\hline

\textbf{Pertemuan 4} &
\textit{Data Preprocessing} 2: Encoding kategorikal, \textit{Scaling/Normalization}, \textit{Feature Engineering} dasar. &
\textbf{Sub-CPMK 4:} Mahasiswa mampu melakukan transformasi data dan rekayasa fitur (\textit{feature engineering}) agar data optimal untuk proses pemodelan. &
\textbf{CPMK2} (fokus eksplorasi data dan pemrograman). \\
\hline

\textbf{Pertemuan 5} &
\textit{Exploratory Data Analysis} (EDA), \textit{Univariate \& Bivariate Analysis}, \textit{Data Visualization} (Matplotlib, Seaborn, dll.). &
\textbf{Sub-CPMK 5:} Mahasiswa mampu merancang dan mengeksekusi Analisis Data Eksploratif (EDA) serta memvisualisasikannya untuk menemukan pola tersembunyi. &
\textbf{CPMK2} (fokus eksplorasi data dan pemrograman). \\
\hline

\textbf{Pertemuan 6} &
Pengantar \textit{Machine Learning}, Analisis Statistik Dasar, Taksonomi Algoritma. &
\textbf{Sub-CPMK 6:} Mahasiswa mampu merumuskan hipotesis dan memilih algoritma yang relevan berdasarkan jenis masalah (\textit{Supervised} vs \textit{Unsupervised Learning}). &
\textbf{CPMK3} (fokus perancangan dan pemilihan algoritma). \\
\hline

\textbf{Pertemuan 7} &
\textit{Supervised Learning: Simple \& Multiple Linear Regression}, Asumsi Regresi. &
\textbf{Sub-CPMK 7:} Mahasiswa mampu merancang dan mengimplementasikan algoritma regresi untuk menyelesaikan permasalahan prediksi kuantitatif pada organisasi. &
\textbf{CPMK3} (fokus perancangan dan pemilihan algoritma). \\
\hline

\textbf{Pertemuan 8} (UTS) &
Evaluasi Pemahaman Konsep, EDA, dan Prapemrosesan Data (Kasus Bisnis Tahap Awal). &
--- &
Menguji pemahaman mahasiswa terhadap \textbf{CPMK1} dan \textbf{CPMK2}. \\
\hline

\textbf{Pertemuan 9} &
\textit{Supervised Learning: Logistic Regression, K-Nearest Neighbors} (KNN), \textit{Decision Tree}. &
\textbf{Sub-CPMK 8:} Mahasiswa mampu mengimplementasikan algoritma klasifikasi dasar untuk memprediksi kelas atau kategori dari suatu permasalahan bisnis. &
\textbf{CPMK3} (fokus perancangan dan pemilihan algoritma). \\
\hline

\textbf{Pertemuan 10} &
\textit{Supervised Learning --- Ensemble: Random Forest, Support Vector Machine} (SVM), \textit{Gradient Boosting}. &
\textbf{Sub-CPMK 9:} Mahasiswa mampu menerapkan algoritma klasifikasi tingkat lanjut (\textit{ensemble}) untuk meningkatkan keandalan dan akurasi prediksi model. &
\textbf{CPMK3} (fokus perancangan dan pemilihan algoritma). \\
\hline

\textbf{Pertemuan 11} &
\textit{Unsupervised Learning: K-Means Clustering, Hierarchical Clustering}, Analisis Silhouette. &
\textbf{Sub-CPMK 10:} Mahasiswa mampu merancang algoritma klastering untuk mengidentifikasi segmentasi data yang tidak berlabel (misal: segmentasi pelanggan). &
\textbf{CPMK3} (fokus perancangan dan pemilihan algoritma). \\
\hline

\textbf{Pertemuan 12} &
Evaluasi Model: \textit{Confusion Matrix, Accuracy, Precision, Recall, F1-Score}, RMSE, ROC-AUC. &
\textbf{Sub-CPMK 11:} Mahasiswa mampu menganalisis dan menghitung metrik evaluasi kinerja model/algoritma yang telah dibangun secara presisi. &
\textbf{CPMK4} (fokus evaluasi dan penyampaian solusi). \\
\hline

\textbf{Pertemuan 13} &
\textit{Cross-Validation} (K-Fold), \textit{Hyperparameter Tuning} (\textit{Grid Search}, \textit{Random Search}). &
\textbf{Sub-CPMK 12:} Mahasiswa mampu mengoptimalkan kinerja algoritma melalui teknik validasi dan penyesuaian parameter (\textit{hyperparameter tuning}). &
\textbf{CPMK3} dan \textbf{CPMK4}. \\
\hline

\textbf{Pertemuan 14} &
\textit{Data Storytelling}, Penyusunan Laporan Bisnis (\textit{Executive Summary}), Rekomendasi Aksi. &
\textbf{Sub-CPMK 13:} Mahasiswa mampu mengkomunikasikan hasil evaluasi teknis algoritma menjadi narasi dan visualisasi (\textit{insight}) yang relevan untuk pengambilan keputusan bisnis. &
\textbf{CPMK4} (fokus evaluasi dan penyampaian solusi). \\
\hline

\textbf{Pertemuan 15} &
Presentasi Proyek Akhir (\textit{Showcase}): Integrasi \textit{Business Understanding, Data Prep, Modeling}, Evaluasi, dan \textit{Business Insight}. &
\textbf{Sub-CPMK 14:} Mahasiswa mampu merancang dan mengeksekusi proyek data science secara utuh (\textit{end-to-end pipeline}) sebagai solusi strategis untuk studi kasus nyata. &
Integrasi utuh dari \textbf{CPMK1}, \textbf{CPMK2}, \textbf{CPMK3}, dan \textbf{CPMK4}. \\
\hline

\textbf{Pertemuan 16} (UAS) &
Evaluasi Pemodelan, Evaluasi Algoritma, dan Penyampaian \textit{Insight} Bisnis (Kasus Bisnis \textit{End-to-End}). &
--- &
Menguji kompetensi utama lanjutan pada \textbf{CPMK3} dan \textbf{CPMK4}. \\
\hline
\end{longtable}
}

\vspace{0.5cm}

% ============================================================
% 6. METODE PEMBELAJARAN
% ============================================================
\section*{6. Metode Pembelajaran}

Strategi atau pendekatan pembelajaran yang dipilih sesuai OBE yang menekankan aktivitas mahasiswa (komposisi 60\% teori, 40\% praktek):

\begin{itemize}[leftmargin=*]
  \item \textbf{Ceramah Interaktif:} Penjelasan konsep data science, statistik, dan machine learning dengan diskusi tanya jawab
  \item \textbf{Problem-Based Learning (PBL):} Mahasiswa menyelesaikan permasalahan analisis data nyata menggunakan Python (NumPy, Pandas, scikit-learn)
  \item \textbf{Project-Based Learning:} Proyek kelompok analisis data end-to-end (data, preprocessing, model, interpretasi) dengan Python
  \item \textbf{Praktikum Terbimbing:} Latihan coding dengan Python (Jupyter/IDE) dengan bimbingan asisten; fokus pada Pandas, Matplotlib/Seaborn, scikit-learn
  \item \textbf{Flipped Classroom:} Mahasiswa mempelajari materi dan tutorial Python sebelum kelas, diskusi mendalam dan praktik di kelas
  \item \textbf{Studi Kasus:} Analisis dataset terbuka (UCI, Kaggle, data.go.id, BPS) dengan Python
  \item \textbf{Peer Review:} Mahasiswa melakukan review notebook/kode Python rekan dan memberikan feedback konstruktif
\end{itemize}

\vspace{0.5cm}

% ============================================================
% 7. PENGALAMAN BELAJAR MAHASISWA
% ============================================================
\section*{7. Pengalaman Belajar Mahasiswa}

Deskripsi tugas, aktivitas, atau pengalaman belajar yang mendukung pencapaian Sub-CPMK:

\begin{itemize}[leftmargin=*]
  \item Melakukan instalasi lingkungan Python (Anaconda/Jupyter atau IDE) dan menjalankan script dasar
  \item Membaca dan memuat dataset dari CSV/Excel menggunakan Pandas
  \item Melakukan pembersihan data (missing, outlier) dan statistik deskriptif dengan Python
  \item Membuat visualisasi (plot) untuk eksplorasi data menggunakan Matplotlib/Seaborn
  \item Menerapkan model klasifikasi atau regresi pada dataset terbuka dengan scikit-learn
  \item Mengevaluasi performa model dan menginterpretasikan metrik (akurasi, RMSE, confusion matrix)
  \item Berkolaborasi dalam tim untuk proyek analisis data lengkap (dari data hingga presentasi)
  \item Mempresentasikan hasil analisis dan insight di depan kelas
  \item Melakukan peer review terhadap notebook/kode Python rekan
  \item Merefleksikan etika data dan privasi dalam penggunaan dataset
\end{itemize}

\vspace{0.5cm}

% ============================================================
% 8. KRITERIA, INDIKATOR, DAN BOBOT PENILAIAN
% ============================================================
\section*{8. Kriteria, Indikator, dan Bobot Penilaian}

Teknik/alat asesmen dipetakan ke Sub-CPMK/CPMK dengan bobot yang jelas. Komposisi: \textbf{60\% teori}, \textbf{40\% praktek}.

\begin{longtable}{|>{\raggedright\arraybackslash}p{2.5cm}|>{\raggedright\arraybackslash}p{4cm}|>{\raggedright\arraybackslash}p{5.5cm}|c|}
\hline
\textbf{Komponen} & \textbf{Teknik Asesmen} & \textbf{Indikator/CPMK} & \textbf{Bobot (\%)} \\
\hline
\endfirsthead
\hline
\textbf{Komponen} & \textbf{Teknik Asesmen} & \textbf{Indikator/CPMK} & \textbf{Bobot (\%)} \\
\hline
\endhead
\hline
\endfoot

Kuis & Tertulis (teori) & Sub-CPMK 1, 6, 11, 12 & 10 \\
\hline
UTS & Ujian tertulis (teori \& pemahaman konsep) & CPMK1, CPMK2 & 25 \\
\hline
UAS & Ujian tertulis komprehensif (teori) & CPMK1 s.d. CPMK4 & 25 \\
\hline
Tugas/Praktikum & Tugas coding \& laporan (Python/R) & Sub-CPMK 2, 3, 4, 5, 7, 8 & 15 \\
\hline
Proyek Kelompok & Analisis data \& presentasi (Python/R) & CPMK2, CPMK3, CPMK4; Sub-CPMK 14 & 25 \\
\hline
\textbf{Total} & & & \textbf{100} \\
\hline
\end{longtable}

\textbf{Rincian komposisi:} Teori (Kuis + UTS + UAS) = 60\%; Praktek (Tugas/Praktikum + Proyek Kelompok) = 40\%.

\textbf{Kriteria Penilaian:}
\begin{itemize}[leftmargin=*]
  \item A (85--100): Menguasai semua CPMK dengan sangat baik, mampu menerapkan dalam kasus kompleks
  \item B (70--84): Menguasai sebagian besar CPMK dengan baik
  \item C (60--69): Menguasai CPMK dasar dengan cukup
  \item D (50--59): Menguasai sebagian kecil CPMK
  \item E (<50): Belum menguasai CPMK yang ditetapkan
\end{itemize}

\vspace{0.5cm}

% ============================================================
% 9. EVALUASI DAN REFLEKSI PEMBELAJARAN (OPSIONAL)
% ============================================================
\section*{9. Evaluasi dan Refleksi Pembelajaran}

Penilaian sumatif/formatif untuk memantau ketercapaian outcome secara menyeluruh:

\begin{itemize}[leftmargin=*]
  \item \textbf{Evaluasi Formatif:} Kuis mingguan, latihan coding Python, peer review untuk memberikan feedback berkelanjutan
  \item \textbf{Evaluasi Sumatif:} UTS dan UAS untuk mengukur pencapaian CPMK secara komprehensif
  \item \textbf{Refleksi Mahasiswa:} Jurnal belajar atau refleksi mingguan terhadap pemahaman materi data science dan Python
  \item \textbf{Evaluasi Dosen:} Survey kepuasan mahasiswa di tengah dan akhir semester
  \item \textbf{Continuous Improvement:} Analisis hasil penilaian untuk perbaikan RPS di semester berikutnya
\end{itemize}

\vspace{0.5cm}

% ============================================================
% 10. DAFTAR REFERENSI (sumber terbuka dari internet)
% ============================================================
\section*{10. Daftar Referensi}

Sumber belajar terbuka (\textit{open access}) dari internet; diselaraskan dengan berkas \texttt{sumber\_materi.md} (pemetaan ke CPL, CPMK, Sub-CPMK, dan materi per pertemuan). Alat utama praktik: Python.

\begin{enumerate}[leftmargin=*]
  \item IBM. \textit{SPSS Modeler CRISP-DM Guide} (PDF). \url{https://public.dhe.ibm.com/software/analytics/spss/documentation/modeler/18.0/en/ModelerCRISPDM.pdf}.

  \item IBM. \textit{CRISP-DM Help Overview} (dokumentasi web). \url{https://www.ibm.com/docs/en/spss-modeler/saas?topic=dm-crisp-help-overview}.

  \item Requests contributors. \textit{Requests: HTTP for Humans} (dokumentasi). \url{https://docs.python-requests.org/en/stable/}. \textit{Quickstart}. \url{https://docs.python-requests.org/en/stable/user/quickstart/}.

  \item Richardson, L. \textit{Beautiful Soup 4 Documentation}. \url{https://www.crummy.com/software/BeautifulSoup/bs4/doc/}.

  \item pandas development team. \textit{User Guide: IO tools} (CSV, Excel, basis data). \url{https://pandas.pydata.org/docs/user_guide/io.html}.

  \item pandas development team. \textit{User Guide: Missing data}. \url{https://pandas.pydata.org/docs/user_guide/missing_data.html}.

  \item scikit-learn developers. \textit{Preprocessing data}. \url{https://scikit-learn.org/stable/modules/preprocessing.html}.

  \item Matplotlib development team. \textit{Matplotlib Tutorials}. \url{https://matplotlib.org/stable/tutorials/index.html}.

  \item Seaborn developers. \textit{Seaborn tutorial}. \url{https://seaborn.pydata.org/tutorial.html}.

  \item scikit-learn developers. \textit{Supervised learning}. \url{https://scikit-learn.org/stable/supervised_learning.html}.

  \item scikit-learn developers. \textit{Unsupervised learning}. \url{https://scikit-learn.org/stable/unsupervised_learning.html}.

  \item scikit-learn developers. Contoh \textit{Selecting the number of clusters with silhouette analysis on KMeans}. \url{https://scikit-learn.org/stable/auto_examples/cluster/plot_kmeans_silhouette_analysis.html}. Fungsi \texttt{sklearn.metrics.silhouette\_score}. \url{https://scikit-learn.org/stable/modules/generated/sklearn.metrics.silhouette_score.html}.

  \item scikit-learn developers. \textit{Metrics and scoring: quantifying the quality of predictions}. \url{https://scikit-learn.org/stable/model_evaluation.html}.

  \item scikit-learn developers. \textit{Cross-validation: evaluating estimator performance}. \url{https://scikit-learn.org/stable/modules/cross_validation.html}. \textit{Tuning the hyper-parameters of an estimator}. \url{https://scikit-learn.org/stable/modules/grid_search.html}.

  \item James, G., Witten, D., Hastie, T., Tibshirani, R., dkk. \textit{An Introduction to Statistical Learning} (situs resmi; edisi R/Python, PDF terbuka). \url{https://www.statlearning.com/}.

  \item VanderPlas, J. \textit{Python Data Science Handbook} (online). \url{https://jakevdp.github.io/PythonDataScienceHandbook/}.

  \item Harvard University. \textit{Data Science: Visualization} (Professional \& Lifelong Learning). \url{https://pll.harvard.edu/course/data-science-visualization}.

  \item Pemerintah Indonesia. \textit{Data Terbuka Indonesia} (\texttt{data.go.id}). \url{https://data.go.id/}.
\end{enumerate}

\vspace{1cm}

\begin{flushright}
\begin{tabular}{c}
Disusun oleh,\\[2cm]
\textbf{[Nama Dosen, M.Kom.]}\\
NIP. [NIP]
\end{tabular}
\end{flushright}

\end{document}
